\documentclass[twocolumn]{article}
\usepackage[english]{babel}

\usepackage[letterpaper,top=2cm,bottom=2cm,left=3cm,right=3cm,marginparwidth=1.75cm]{geometry}
\usepackage{amsmath}
\usepackage{graphicx}
\usepackage[colorlinks=true, allcolors=red]{hyperref}
\usepackage{booktabs}
\usepackage{array}
\usepackage{multirow}
\usepackage{float}
\usepackage{mwe}             
\usepackage{algorithm}       
\usepackage{algorithmic}     
\usepackage{natbib}          

\title{\textbf{Compact Image Classification Pipeline}}
\author{A. Student\\
Department of Computer Science and Engineering\\
Rajshahi University of Engineering \& Technology\\
student@cse.ruet.ac.bd
}
\date{}

\begin{document}
\maketitle

\begin{abstract}
This one-page document summarizes a compact image classification pipeline, illustrating its preprocessing outputs, performance table, processing flow, and training procedure. The layout follows a two-column journal format with Harvard-style author-year referencing, a $2\times2$ figure panel confined to a single column, and a full-width schematic of the inference flow.\\\\
\textbf{Keywords:} image classification, convolutional neural networks, preprocessing, model pruning.
\end{abstract}

\section{Preprocessing Examples}
Each input image passes through four deterministic stages before it reaches the backbone network. Figure~\ref{fig:panel} shows the intermediate output of every stage for one representative sample, arranged as a $2\times2$ panel inside a single column.

\begin{figure}[H]
    \centering
    \begin{minipage}{0.23\textwidth}
        \includegraphics[width=\linewidth]{example-image-a}
    \end{minipage}
    \begin{minipage}{0.23\textwidth}
        \includegraphics[width=\linewidth]{example-image-b}
    \end{minipage}\\[2pt]
    \begin{minipage}{0.23\textwidth}
        \includegraphics[width=\linewidth]{example-image-c}
    \end{minipage}
    \begin{minipage}{0.23\textwidth}
        \includegraphics[width=\linewidth]{example-image-a}
    \end{minipage}
    \caption{Preprocessing stages arranged as a $2\times2$ panel.}
    \label{fig:panel}
\end{figure}

\section{Results}
Table~\ref{tab:basic} reports top-1 accuracy for three backbone variants, evaluated under an identical training budget. Pruning the residual backbone costs 1.5 accuracy points while removing roughly a third of its parameters.

\begin{table}[H]
    \centering
    \caption{Basic table}
    \label{tab:basic}
    \begin{tabular}{lcc}
        \toprule
        Backbone & Variant & Accuracy (\%) \\
        \midrule
        Resnet & Base & 91.2 \\
        Resnet & Pruned & 89.7 \\
        Mobilenet & Base & 88.4 \\
        \bottomrule
    \end{tabular}
\end{table}

\section{Preprocessing Examples}
Each input image passess through four deterministic stages before it reaches the backbone network. Figure 2 
% to wrok with figure 2 thing label
shows the intermediate output of every stage for one representative sample, arranged as a $2\times2$ panel inside a single column.


\section{Training Procedure}
Algorithm 1 lists the optimisation loop. Weights are updated with plain stochastic gradient descent, following the general approach of Lecun \textit{et at.} (1998) and Vaswani {et at.} (2017).\\




\begin{algorithm}[H]
    \caption{Model Training}
    \label{alg:sgd}
    \begin{algorithmic}
        \STATE Input: Training Set D, epochs E, learning rate n
        \STATE Output: Trained weights $\theta$
        \FOR{each epoch}
            \FOR{each mini-batch $(x, y)$}
                \STATE Compute loss $L(w; x, y)$
                \STATE $w \gets w - \eta \nabla L(w)$
            \ENDFOR
        \ENDFOR
    \end{algorithmic}
\end{algorithm}

\bibliographystyle{plainnat}
\begin{thebibliography}{9}
    \bibitem[LeCun et al., 1998]{lecun1998} LeCun, Y., Bottou, L., Bengio, Y. and Haffner, P. (1998) 'Gradient-based learning applied to document recognition', \textit{Proceedings of the IEEE}, 86(11), pp. 2278--2324.
    \bibitem[Vaswani et al., 2017]{vaswani2017} Vaswani, A., Shazeer, N., Parmar, N., Uszkoreit, J., Jones, L., Gomez, A.N., Kaiser, L. and Polosukhin, I. (2017) 'Attention is all you need', \textit{Advances in Neural Information Processing Systems}, 30, pp. 5998--6008.
\end{thebibliography}

\end{document}


@article{lecun1998,
    author  = {LeCun, Y. and Bottou, L. and Bengio, Y. and Haffner, P.},
    title   = {Gradient-based learning applied to document recognition},
    journal = {Proceedings of the IEEE},
    year    = {1998},
    volume  = {86},
    number  = {11},
    pages   = {2278--2324}
}
@article{vaswani2017,
    author  = {Vaswani, A. and Shazeer, N. and Parmar, N. and Uszkoreit, J. and Jones, L. and Gomez, A. N. and Kaiser, L. and Polosukhin, I.},
    title   = {Attention is all you need},
    journal = {Advances in Neural Information Processing Systems},
    year    = {2017},
    volume  = {30},
    pages   = {5998--6008}
}